\documentclass[reqno, answers]{exam}

\usepackage{amssymb, amsmath, amsthm, stmaryrd}
\usepackage{fullpage, parskip}
\usepackage{enumitem}
\usepackage{hyperref}

\title{Introduction to Computational Math Spring 2026 \\ Homework 11}
\date{Due: Friday, Apr 17, 2026, 11:59 PM}
\author{}

\begin{document}
\maketitle
\thispagestyle{empty}

Textbook = Driscoll and Braun (\url{https://fncbook.com/})

\begin{enumerate}

  \item Textbook Exercise 6.4.1a
        \begin{solution}
          $u_1 = \boxed{1.92}$ at $t_1 = 0.2$.

          Exact: $\hat{u}(0.2) = 2e^{-0.04} \approx 1.9216$; error $\approx 0.0016$.
        \end{solution}
  \item Textbook Exercise 6.4.2a
        \begin{solution}
          $u_1 = \boxed{1.92}$ at $t_1 = 0.2$.

          Exact: $\hat{u}(0.2) = 2e^{-0.04} \approx 1.9216$; error $\approx 0.0016$.
        \end{solution}
  \item Textbook Exercise 6.4.6a
        \begin{solution}
          $k_1 = 0$, $k_2 = -0.4$, $k_3 = -0.392$, $k_4 = -0.76864$.

          $u_1 = \boxed{1.9216}$ at $t_1 = 0.2$ (essentially exact; error $\approx 0$).
        \end{solution}
  \item Textbook Exercise 6.4.11
        \begin{solution}
          \textbf{(a)}
          \[
            \frac{u_{i+1}}{u_i} = 1 + ch + \frac{(ch)^2}{2}
          \]

          \textbf{(b)} When $ch = -3$: $\dfrac{u_{i+1}}{u_i} = 1 - 3 + \dfrac{9}{2} = 2.5$, so $u_i = (2.5)^i \to \infty$.
          The exact solution is $\hat{u}(t) = e^{ct} \to 0$ since $c < 0$.
          The method is unstable for $ch = -3$.
        \end{solution}


  \item It is possible to generalize Euler's method using any finite difference approximation to $u''(t)$.
        \begin{parts}
          \part Let $k \in (0, 1)$ be a constant. Construct a FD approximation for $u''(t)$ using $u'(t)$ and $u'(t+kh)$.
          \part Derive a new one-step method for solving $u' = f(t,u)$.
          \part What is the order of accuracy of this method?
          \part Write the Butcher tableau for this method.
        \end{parts}

        \begin{solution}
          \begin{parts}
            \part $\displaystyle u''(t) = \frac{u'(t+kh) - u'(t)}{kh} + O(h)$
            \part
            \begin{align*}
              k_1     & = hf(t_i, u_i) \\
              k_2     & = hf(t_i + kh,\, u_i + k\, k_1) \\
              u_{i+1} & = u_i + \left(1 - \frac{1}{2k}\right) k_1 + \frac{1}{2k} k_2
            \end{align*}
            \part Order 2.
            \part
            \[
              \begin{array}{c|cc}
                0 &                  &              \\[0.5em]
                k & k                &              \\[0.5em]
                \hline
                  & 1 - \frac{1}{2k} & \frac{1}{2k} \\
              \end{array}
            \]
          \end{parts}
        \end{solution}
  \item The classic Runge-Kutta method (called the 3/8th rule) uses a slightly different Butcher tableau than the one we discussed in class:
        \[
          \begin{array}{c|cccc}
            0   &      &     &           \\[0.5em]
            1/3 & 1/3  &     &           \\[0.5em]
            2/3 & -1/3 & 1   &           \\[0.5em]
            1   & 1    & -1  & 1         \\[0.5em]
            \hline
                & 1/8  & 3/8 & 3/8 & 1/8 \\
          \end{array}
        \]
        Note that any \emph{missing} entries in the Butcher tableau are assumed to be zero.
        \begin{parts}
          \part Explicitly write out the update formula for this method.
          \part What integration method does this correspond to when applied to the integral form of the IVP? (This is in fact another numerical integration method called Simpson's 3/8th rule.)
        \end{parts}
        \begin{solution}
          \begin{parts}
            \part
            \begin{align*}
              k_1     & = hf(t_i, u_i) \\
              k_2     & = hf\!\left(t_i + \tfrac{h}{3},\, u_i + \tfrac{1}{3}k_1\right) \\
              k_3     & = hf\!\left(t_i + \tfrac{2h}{3},\, u_i - \tfrac{1}{3}k_1 + k_2\right) \\
              k_4     & = hf\!\left(t_i + h,\, u_i + k_1 - k_2 + k_3\right) \\
              u_{i+1} & = u_i + \tfrac{1}{8}k_1 + \tfrac{3}{8}k_2 + \tfrac{3}{8}k_3 + \tfrac{1}{8}k_4
            \end{align*}
            \part Simpson's 3/8th rule:
            \[
              \int_{t_i}^{t_{i+1}} u'(t)\,dt \approx \frac{h}{8}\left[f(t_i) + 3f\!\left(t_i+\tfrac{h}{3}\right) + 3f\!\left(t_i+\tfrac{2h}{3}\right) + f(t_{i+1})\right]
            \]
          \end{parts}
        \end{solution}

  \item Jupyter Notebook

\end{enumerate}

\section*{Quiz Information}

The quiz will be held on {\bf Tuesday, April 21, 2026} during the discussion section.

Quiz questions will be modifications of the {\it homework problems and material discussed in class}.

\textbf{Topics covered:}

Chapters 6.4

\begin{itemize}
  \item Runge-Kutta methods for solving IVPs
\end{itemize}

You should also review Chapter 5.4 on finite difference approximations for derivatives.
You should memorize the basic Euler's method. I'll provide a Butcher tableau for any other methods that I ask about on the quiz.



\end{document}